% !TEX TS-program = pdflatex
% !TEX encoding = UTF-8 Unicode

% This file is a template using the "beamer" package to create slides for a talk or presentation
% - Giving a talk on some subject.
% - The talk is between 15min and 45min long.
% - Style is ornate.

% MODIFIED by Jonathan Kew, 2008-07-06
% The header comments and encoding in this file were modified for inclusion with TeXworks.
% The content is otherwise unchanged from the original distributed with the beamer package.

\documentclass{beamer}
\setbeamercovered{invisible}



% Copyright 2004 by Till Tantau <tantau@users.sourceforge.net>.
%
% In principle, this file can be redistributed and/or modified under
% the terms of the GNU Public License, version 2.
%
% However, this file is supposed to be a template to be modified
% for your own needs. For this reason, if you use this file as a
% template and not specifically distribute it as part of a another
% package/program, I grant the extra permission to freely copy and
% modify this file as you see fit and even to delete this copyright
% notice. 


\mode<presentation>
{
  \usetheme{Warsaw}
  % or ...

  % or whatever (possibly just delete it)
}


\usepackage[english]{babel}
% or whatever

\usepackage[utf8]{inputenc}
% or whatever

\usepackage{times}
\usepackage[T1]{fontenc}
% Or whatever. Note that the encoding and the font should match. If T1
% does not look nice, try deleting the line with the fontenc.

%%% MATH RELATED

\input{macros.tex}

\title{MATH 350-2 Advanced Calculus} 
\subtitle
{} % (optional)

\author[W.R. Casper] % (optional, use only with lots of authors)
{W.R. Casper}
% - Use the \inst{?} command only if the authors have different
%   affiliation.

\institute[California State University Fullerton] % (optional, but mostly needed)
{
  Department of Mathematics\\
  California State University Fullerton}
% - Use the \inst command only if there are several affiliations.
% - Keep it simple, no one is interested in your street address.

\subject{Talks}
% This is only inserted into the PDF information catalog. Can be left
% out. 



% If you have a file called "university-logo-filename.xxx", where xxx
% is a graphic format that can be processed by latex or pdflatex,
% resp., then you can add a logo as follows:

% \pgfdeclareimage[height=0.5cm]{university-logo}{university-logo-filename}
% \logo{\pgfuseimage{university-logo}}



% Delete this, if you do not want the table of contents to pop up at
% the beginning of each subsection:
\AtBeginSubsection[]
{
  \begin{frame}<beamer>{Outline}
    \tableofcontents[currentsection,currentsubsection]
  \end{frame}
}


% If you wish to uncover everything in a step-wise fashion, uncomment
% the following command: 

%\beamerdefaultoverlayspecification{<+->}


\begin{document}

\begin{frame}
  \titlepage
\end{frame}

\begin{frame}{Outline}
  \tableofcontents
  % You might wish to add the option [pausesections]
\end{frame}

% Since this a solution template for a generic talk, very little can
% be said about how it should be structured. However, the talk length
% of between 15min and 45min and the theme suggest that you stick to
% the following rules:  

% - Exactly two or three sections (other than the summary).
% - At *most* three subsections per section.
% - Talk about 30s to 2min per frame. So there should be between about
%   15 and 30 frames, all told.

\section{Real Analysis Lecture 23}

\subsection{Differentiation}

\begin{frame}{Differentiability}
\begin{defn}
A real-valued function $f(x)$ defined on an open interval $(a,b)$ is \textbf{differentiable} at $c\in (a,b)$ if
$$\lim_{x\rightarrow c}\frac{f(x)-f(c)}{x-c}\ \ \text{exists}.$$
In this case the limit is called the derivative of $f$ at $c$, and denote by $f'(c)$.
\end{defn}
\end{frame}

\begin{frame}{Chord function}
\begin{thm}[Helper Theorem]
Let $f$ be a real-valued function on $(a,b)$.  Then $f$ is differentiable at $c\in (a,b)$ if and only if there exists a function $f^*$ on $(a,b)$ which is continuous at $c$ and satisfies
$$f(x)-f(c) = (x-c)f^*(x).$$
\end{thm}
\pause
The function $f^*(x)$ is the \textbf{chord slope function}.
\end{frame}

\begin{frame}{Chord function}
\begin{proof}
\pause
Suppose $f(x)$ is differentiable at $c$.
\pause
Define
$$f^*(x) = \left\lbrace\begin{array}{cc}
\frac{f(x)-f(c)}{x-c}, & x\neq c\\
f'(c), & x= c
\end{array}\right.$$
\pause
Then $\lim_{x\rightarrow c} f^*(x) = f'(c)$, so $f^*$ is continuous at $c$.\\
\pause
Conversely, suppose $f^*$ exists and is continuous at $c$.\\
\pause
Then
$$\lim_{x\rightarrow c} \frac{f(x)-f(c)}{x-c} = \lim_{x\rightarrow c} f^*(x) = f^*(c).$$
\pause
In particular $f$ is differentiable at $c$ and $f'(x) = f^*(c)$.
\end{proof}
\end{frame}

\begin{frame}{Differentiable implies continuous}
\begin{thm}
Suppose $f(x)$ is a real valued function on $(a,b)$ and is differentiable at $c\in (a,b)$.\\
Then $f$ is continuous at $c$.
\end{thm}
\pause
\begin{proof}
\pause
By the Helper Theorem, there exists a function $f^*$ on $(a,b)$ which is continuous at $c$ with
$$f(x) = (x-c)f^*(x).$$
\pause
Since $(x-c)$ and $f^*$ are both continuous at $c$, so is $f(x)$.
\end{proof}
\end{frame}

\begin{frame}{Additivity of derivatives}
\begin{thm}
Let $f,g$ be real functions on $(a,b)$ which are differentiable at $c\in (a,b)$.\\
$$(f\pm g)'(c) = f'(c) \pm g'(c)$$
\end{thm}
\end{frame}

\begin{frame}{Additivity of derivatives}
\begin{proof}
\pause
By the Helper Theorem
$$f(x)-f(c) = (x-c)f^*(x),\quad g(x)-g(c) = (x-c)g^*(x)$$
\pause
Therefore
$$(f\pm g)(x) - (f\pm g)(c) = (x-c)(f^*(x) \pm g^*(x)).$$
\pause
Since $f^*$ and $g^*$ are continuous at $c$, so is $f^*\pm g^*$.\\
\pause
The Helper Theorem implies that $g\pm f$ is differentiable and
$$(f\pm g)'(c) = f'(c) \pm g'(c)$$
\end{proof}
\end{frame}

\begin{frame}{Product rule}
\begin{thm}
Let $f,g$ be real functions on $(a,b)$ which are differentiable at $c\in (a,b)$.\\
$$(fg)'(c) = f'(c)g(c) + f(c)g'(c)$$
\end{thm}
\end{frame}
\begin{frame}{Product rule}
\begin{proof}
\pause
By the Helper Theorem
$$f(x)-f(c) = (x-c)f^*(x),\quad g(x)-g(c) = (x-c)g^*(x)$$
\pause
Therefore
\begin{align*}
(fg)(x)-(fg)(c)
  & = f(x)g(x) - f(c)g(c)\\
  & = f(x)g(x) - f(x)g(c) + f(x)g(c)-f(c)g(c)\\
  & = f(x)(g(x) - g(c)) + (f(x)-f(c))g(c)\\
  & = (x-c)[f(x)g^*(x) + f^*(x)g(c)]
\end{align*}
\pause
Helper Theorem says:
$$(fg)'(c) = f(c)g^*(c) + f^*(c)g(c) = f(c)g'(c) + f'(c)g(c).$$
\end{proof}
\end{frame}

\begin{frame}{Quotient rule}
\begin{thm}
Let $f,g$ be real functions on $(a,b)$ which are differentiable at $c\in (a,b)$ with $g(c)\neq 0$.\\
$$(f/g)'(c) = \frac{f'(c)g(c) - f(c)g'(c)}{g(c)^2}$$
\end{thm}
\end{frame}
\begin{frame}{Quotient rule}\mbox{}
{\tiny
By the Helper Theorem
$$f(x)-f(c) = (x-c)f^*(x),\quad g(x)-g(c) = (x-c)g^*(x)$$
\pause
\begin{align*}
(f/g)(x) - (f/g)(c)
  & = \frac{f(x)g(c)-f(c)g(x)}{g(x)g(c)}\\
  & = \frac{f(x)g(c)-f(c)g(c)+f(c)g(c)-f(c)g(x)}{g(x)g(c)}\\
  & = \frac{(f(x)-f(c))g(c)-f(c)(g(x)-g(c))}{g(x)g(c)}\\
  & = (x-c)\frac{f^*(x)g(c)-f(c)g^*(x)}{g(x)g(c)}\\
\end{align*}
\pause
Helper Theorem says:
$$(f/g)'(c) = \frac{f^*(c)g(c)-f(c)g^*(c)}{g(c)^2} = \frac{f'(c)g(c)-f(c)g'(c)}{g(c)^2}.$$
}
\end{frame}

\begin{frame}{Chain rule}
\begin{thm}
Let $f$ be a real function on an interval $I = (a,b)$, and let $g$ be a real function on $f(I)$, and consider the composition
$$h = g\circ f: (a,b)\rightarrow\mathbb{R}.$$
Suppose $c\in (a,b)$ and $f(c)$ is an interior point of $f(I)$.
If $f$ is differentiable at $c$ and $g$ is differentiable at $f(c)$, then $h$ is differentiable at $c$ and
$$h'(c) = (g\circ f)'(c) = g'(f(c))f'(c).$$
\end{thm}
\end{frame}

\subsection{Derivatives and Optima}

\begin{frame}{Local maxima and minima}
\begin{defn}
Let $(M,d)$ be a metric space and $f$ be a real-valued function on a subset $S\subseteq M$.
\pause
Then $f$ has an \textbf{absolute (or global) maximum at $c\in S$} if
$$f(x)\leq f(c)\quad\text{for all}\ x\in S.$$
\pause
It has \textbf{local maximum at $a\in S$} if there exists an $r>0$ such that
$$f(x)\leq f(c)\quad\text{for all}\ x\in B_M(c;r)\cap S.$$
\pause
Absolute and local minima are defined similarly.
\end{defn}
\end{frame}

\begin{frame}{Critical points}
\begin{defn}
Let $f$ be a real-valued function on an open interval $(a,b)$.
\pause
A point $c\in (a,b)$ is called a \textbf{critical point} of $f$ if either $f$ is not differentiable at $c$ or else $f'(c) = 0$.
\end{defn}
\pause
\begin{thm}
Let $f$ be a real-valued function on an open interval $(a,b)$ and suppose that $f$ has a local maximum or local minimum at $c\in (a,b)$.
Then $c$ is a critical point of $f$.
\end{thm}
\end{frame}

\begin{frame}{Critical points}
\begin{proof}
\pause
We will prove this in the case $c$ is a local maximum.  The proof in the case of a local minimum is similar.\\
\pause
If $f$ is not differentiable at $c$, then we're done!\\
\pause
Therefore assume that it is differentiable.\\
\pause
By definition, there exists $r>0$ such that
$$f(x)\leq f(c)\quad\text{for all}\ x\in (c-r,c+r)\cap (a,b).$$
\pause
To obtain a contradiction, suppose that $f'(c)\neq 0$.\\
\pause
Then by the Helper Theorem, there is a function $f^*(x)$ defined on $(a,b)$ and continuous at $c\in (a,b)$ with
$$f(x)-f(c) = (x-c)f^*(x).$$
\end{proof}
\end{frame}

\begin{frame}{Critical points}
\begin{proof}
Moreover, $f^*(c) = f'(c)\neq 0$.\\
\pause
Since $f$ is continuous, we may choose $\delta > 0$ such that
$$|x-c| < \delta\ \ \Rightarrow\ \ |f^*(x) - f^*(c)| < |f^*(c)|/2.$$
\pause
Moreover, we may assume that $\delta < r$.\\
\pause
It follows that $f^*(x) > |f^*(c)|/2$ for all $x\in (a,b)$ with $|x-c|<\delta$.
\pause
However,
$$f(c\pm \delta/2) = f(c) \pm \frac{1}{2}\delta f^*(c+\delta/2),$$
one of these will be larger than $f(c)$.  \pause Contradiction!
\end{proof}
\end{frame}

\subsection{Mean Value Theorem}

\begin{frame}{Rolle's Theorem}
\begin{thm}{Rolle's Theorem}
Suppose that $f(x)$ is a real-valued function on $[a,b]$ which is
\begin{itemize}
\pause
\item continuous on $[a,b]$;
\pause
\item differentiable at each point in $(a,b)$;
\pause
\item and which satisfies $f(a) = f(b)$.
\end{itemize}
\pause
Then there exists $c\in (a,b)$ with $f'(c) = 0$.
\end{thm}
\end{frame}

\begin{frame}{Rolle's Theorem}
\begin{proof}
Suppose otherwise.\\
\pause
By the Extreme Value Theorem, $f(x)$ has an absolute maximum on $[a,b]$, as well as an absolute minimum on $[a,b]$.\\
\pause
If both of these occur at the endpoints, then the absolute max is the absolute min, forcing $f$ to be constant.\\
\pause
Then $f'(c)=0$ for every $c\in (a,b)$.\\
\pause
Otherwise at least one of the points happens inside, so $f(x)$ has a local max or local min at a value $c$ inside $(a,b)$.\\
\pause
The point $c$ is a critical point, and since $f$ is differentiable at $c$, we must have $f'(c)=0$.
\end{proof}
\end{frame}

\begin{frame}{Mean Value Theorem}
\begin{thm}{Rolle's Theorem}
Suppose that $f(x)$ is a real-valued function on $[a,b]$ which is
\begin{itemize}
\pause
\item continuous on $[a,b]$;
\pause
\item and differentiable at each point in $(a,b)$.
\end{itemize}
\pause
Then there exists $c\in (a,b)$ with
$$f'(c) = \frac{f(b)-f(a)}{b-a}.$$
\end{thm}
\end{frame}

\begin{frame}{Mean Value Theorem}
\begin{proof}
Consider the function
$$g(x) = f(x)-\left(\frac{f(b)-f(a)}{b-a}\right)(x-a) + f(a).$$
\pause
Then $g(x)$ is
\begin{itemize}
\pause
\item continuous on $[a,b]$;
\pause
\item differentiable at each point in $(a,b)$;
\pause
\item and satisfies $g(a) = g(b)$.
\end{itemize}
\pause
By Rolle's Theorem, there exists $c\in (a,b)$ with
$$0 = g'(c) = f'(c) -\frac{f(b)-f(a)}{b-a}.$$
\end{proof}
\end{frame}



\end{document}


